\section{Introduction}
In the era of artificial intelligence (AI)\cite{lecun2015deep}, the popularity of AI compilers has increased dramatically~\cite{chen2018tvm, ansel2024pytorch}.
Modern AI compilers, such as Triton\cite{tillet2019triton}, have become popular in bridging the gap between high-level machine learning framework operators (e.g. general matrix multiplications (GEMMs)\cite{nvidia2023cublas}, softmax\cite{rumelhart1986learning}, etc.) and the low-level, target-specific machine code\cite{NVIDIAPTX}. 
The diversity of workloads and the rapid evolution of hardware architectures, particularly in GPUs, pose significant challenges in developing frameworks that are both modular in development and performant ``out of the box'' for users.
Addressing these demands requires a flexible compiler infrastructure, like LLVM\cite{lattner2004llvm}'s Multi-Level Intermediate Representation (MLIR)\cite{lattner2021mlir} with a modular design, and Triton for its customizable high-performance AI operators for 
diverse workloads and backends.

\begin{figure*}
    \centering
    \includegraphics[width=\linewidth]{fig/concept.pdf}
    \vspace{-0.7cm}
    \caption{\R{Concept of the \thiswork{} infrastructure and ecosystem for compiler-centric performance tool. (Left) Overview and comparison of \thiswork{}'s compiler-centric design and prior profiler designs. (Right) Demonstrative examples of novel performance tools facilitated by the compiler-centric design of \thiswork{}.}}
    \vspace{-0.3cm}
    \label{fig:concept}
\end{figure*}

Despite significant advancements, existing compilers often struggle to outperform hand-tuned implementations like cuBLAS\cite{nvidia2023cublas}, rocBLAS\cite{amd2023rocblas}, and CUTLASS\cite{Thakkar_CUTLASS_2023}.
This limitation stems from GPU architectures' rapid evolution and the continuous emergence of novel operator variants. 
For instance, Nvidia’s Hopper architecture incorporates advanced acceleration units such as $5^{\text{th}}$ generation Tensor Cores (TC)\cite{nvidia2018turing} and Tensor Memory Accelerators (TMA)\cite{choquette2023nvidia}, which require sophisticated compiler passes to fully leverage their potential with a dataflow-oriented programming paradigm\cite{choquette2023nvidia}. 
Meanwhile, the latest Flash-Attention-3\cite{shahflashattention} (FA3) kernel, crucial for large language models (LLMs)\cite{brown2020language, llm_for_finance, llm_for_medical}, employs complex tiling and pipelining techniques that traditional compilers cannot effectively handle.

To enhance the compilers with novel compute paradigms and hardware features, profiling tools are critical in developing high-performance AI kernels and optimization passes. 
They are essential for identifying performance bottlenecks, analyzing execution flow, and understanding memory access patterns.
Developers depend on these tools to pinpoint instruction stalls, optimize kernel performance, and refine compiler passes. 
Even for feedback-guided auto-optimization compilers\cite{zheng2020ansor}, the design space consists of manually articulated transformation primitives and requires precise performance insights to guide effective optimization. 
Without effective profiling tools, achieving peak performance on modern GPU platforms becomes a daunting challenge.

However, existing profilers are poorly aligned with the advancements of compiler infrastructure to assist in the development of themselves and AI kernels.
\textbf{Our key observation is that prior profiler designs are isolated from the compiler system and lack key connections to the upper-level framework operations.}
Profiling tools are usually developed as external tools, detached from the intricacies of the compilation process, and fail to provide framework operation-informed, actionable insights tailored to the unique challenges of AI compilers. 
This disconnection severely hinders the ability of operator and AI compiler developers to enhance the performance of ML workloads.


To address this, we propose a novel compiler-centric infrastructure, \thiswork{}, to facilitate the development of customizable and reusable performance tools\footnote{\OURR{The source code is open sourced at \url{https://github.com/triton-lang/triton/tree/main/third_party/proton/dialect}}.}, as illustrated in \Fig{\ref{fig:concept}}.
Our approach establishes a seamless interplay between the compiler and profiler, enabling the exchange of program semantics and profiling metrics. 
\OURR{
Central to our methodology is incorporating multi-level compiler IR instrumentation into the profiling workflow. 
This allows performance tools to be composed as compiler passes with compiler-integrated profiling operations.
This design fundamentally enhances performance tooling in three key ways.
}

\textbf{\ballnumber{1} Programmable tools with IR semantics.} Existing profilers have limited information on the program, restricting their ability to capture performance details tied to program semantics. For example, traditional tools cannot track how software pipelining\cite{huang_alcop_2023} overlaps and evolves across pipeline stages due to missing loop-level information. By integrating IR-level instrumentation, \thiswork{} enables profiling tools to capture these nuanced behaviors effectively.

\textbf{\ballnumber{2} Profile-driven design of compiler passes.} Many optimization techniques, such as autotuning\cite{tunner1, tunner-autotvm}, rely on performance feedback to guide their tuning processes\cite{zheng2020ansor}. Traditional approaches treat compilers and profilers as isolated sub-modules, requiring standalone systems to coordinate optimization and evaluation. With \thiswork{}, profiling passes can directly interact with compiler optimization passes, streamlining the feedback loop and enabling more optimizations.

\textbf{\ballnumber{3} Reusable and extendable performance tools.} By implementing tools native to the compiler IR, \thiswork{} ensures their reusability and portability across frameworks and backends. Profiling tools can operate on shared upper-level representations used by AI frameworks and seamlessly lower to diverse backends, including Nvidia and AMD GPU platforms. Building on the MLIR philosophy of a shared set of infrastructure, we provide not only a powerful set of analysis tools but also a common set of building blocks and infrastructure for a community-driven ecosystem.

To showcase the capabilities of the \thiswork{} infrastructure, we present the first region-based timing tool for GPUs to provide insights into intra-kernel behaviors. 
By leveraging the program's IR semantics, the tool efficiently updates the profiled regional timestamps to recover the inaccuracy caused by instrumentation, ensuring accurate profiling results while minimizing memory overhead. 
This level of precision was previously difficult to achieve because traditional profilers lack access to compiler-generated semantic information, such as loop structures and regional boundaries, making it challenging to attribute execution metrics accurately to high-level behaviors. 
This timing tool plays a crucial role in optimizing modern GPU kernels by enabling precise analysis of fine-grained overlapping and presenting an intuitive timeline visualization to guide performance optimizations.

We conducted an in-depth case study on the FA3 kernel. 
Leveraging the tool's fine-grained profiling results, we identified idle bubble regions in the baseline implementation and extracted key optimization insights. 
Based on these findings, we implemented novel compiler passes to enhance the kernel with an improved overlapping strategy, effectively reducing performance bottlenecks.
In addition, we introduced performance modeling based on the profiling results to evaluate the efficiency of the optimized overlapping design.

This work makes the following key contributions:
\begin{itemize}
    \itemsep0em 
    \item We propose \thiswork{}, a compiler-centric infrastructure enabling customizable performance tools. Integrated into the Triton compiler, it supports AMD and Nvidia platforms, providing essential support for developers optimizing AI compilers and operators (\Sec{\ref{sec:compiler_centric}}).
    \item As a demonstration of the proposed infrastructure, we develop a novel region-based timing tool. This tool exemplifies the compiler-centric approach and offers fine-grained GPU profiling capabilities (\Sec{\ref{sec:tool}}).
    \item We perform an in-depth case study using the proposed timing tool to analyze GPU overlapping. The study reveals valuable insights into intra-kernel overlapping techniques, such as warp specialization (\Sec{\ref{sec:evaluation}}).
\end{itemize}